% GENERATED FILE. Edit canonical lessons, then run npm run generate:books.
% canonical-source-hash: fnv1a64:03bfa32daf123e77
% canonical-lessons: AR-C541-fannan, AR-C541-marid, AR-C541-nahhat, AR-C541-nahata, AR-C541-rassam

\chapter{At the exhibition}
\label{ch:ar-at-the-exhibition}
\hlchaptermodality{541}

\begin{chapteropening}
\textbf{By the end of this chapter:} \emph{I can name an artist, an exhibition, a sculptor and a painter, and say that someone carved a figure.}

Complete the last lesson of chapter 541: I can name an artist, an exhibition, a sculptor and a painter, and say that someone carved a figure.
\end{chapteropening}

\section[\texorpdfstring{\ar{فنان} (fannān)}{فنان (fannān)}]{\ar{فنان} (fannān) — an artist}

\label{lesson:AR-C541-fannan}



\begin{quote}
Before the new one: say the Arabic for passion, then the Arabic for fondness.
\end{quote}

\subsection*{You'll want to know: \ar{فنان}}
\textbf{\ar{فنان}} — \emph{fannān} — \textquotedblleft{}an artist\textquotedblright{}.

\begin{quote}
\ar{رسم} \ar{الفنان} \ar{صورة} \ar{لجدتي}.
\end{quote}

\begin{quote}
\emph{rasama al-fannān ṣūra li-jaddatī.}
\end{quote}

\begin{quote}
The artist painted a picture of my grandmother.
\end{quote}

\begin{cousinweb}
\textbf{\ar{فنان}} (\emph{fannān}) is built on the native Arabic root \emph{f-n-n}, \textquotedblleft{}a sort, a kind; a branch\textquotedblright{} (Lane). Fann, an art, which you know, is from the same root.
\end{cousinweb}

\begin{grammarlens}[title={What you've built}]
One more word for describing what you lived and felt.
\end{grammarlens}

\subsection*{Your turn}
\begin{itemize}
\raggedright
  \item \emph{Say it:} \emph{fannān}
  \item \emph{Say it:} \emph{fannān}, once more
  \item \emph{Say it:} \emph{rasama al-fannān ṣūra li-jaddatī.} — The artist painted a picture of my grandmother.
  \item \emph{Recall:} say the Arabic for passion, then the Arabic for fondness, then say \emph{fannān} again
\end{itemize}

\begin{tcolorbox}[breakable,colback=teal!4,colframe=teal!35!black,title={Before you move on}]
What does \ar{فنان} mean? (\textquotedblleft{}An artist\textquotedblright{}.) Say it once more.
\end{tcolorbox}
\section[\texorpdfstring{\ar{معرض} (maʿriḍ)}{معرض (maʿriḍ)}]{\ar{معرض} (maʿriḍ) — an exhibition}

\label{lesson:AR-C541-marid}



\begin{quote}
Before the new one: say the Arabic for fondness, then the Arabic for an artist.
\end{quote}

\subsection*{You'll want to know: \ar{معرض}}
\textbf{\ar{معرض}} — \emph{maʿriḍ} — \textquotedblleft{}an exhibition\textquotedblright{}.

\begin{quote}
\ar{زرنا} \ar{معرض} \ar{الكتاب} \ar{في} \ar{المدينة}.
\end{quote}

\begin{quote}
\emph{zurnā maʿriḍ al-kitāb fī al-madīna.}
\end{quote}

\begin{quote}
We visited the book fair in the city.
\end{quote}

Plural: \ar{معارض}, \emph{maʿāriḍ}.

\ar{معرض} \ar{الكتاب} is a book fair.

\begin{cousinweb}
\textbf{\ar{معرض}} (\emph{maʿriḍ}) is built on the native Arabic root \emph{ʿ-r-ḍ}, \textquotedblleft{}to appear, to come in the way\textquotedblright{} (Lane). A maʿriḍ is a place where things are shown.
\end{cousinweb}

\begin{grammarlens}[title={What you've built}]
One more word for describing what you lived and felt.
\end{grammarlens}

\subsection*{Your turn}
\begin{itemize}
\raggedright
  \item \emph{Say it:} \emph{maʿriḍ}
  \item \emph{Say it:} \emph{maʿriḍ}, once more
  \item \emph{Say it:} \emph{zurnā maʿriḍ al-kitāb fī al-madīna.} — We visited the book fair in the city.
  \item \emph{Recall:} say the Arabic for fondness, then the Arabic for an artist, then say \emph{maʿriḍ} again
\end{itemize}

\begin{tcolorbox}[breakable,colback=teal!4,colframe=teal!35!black,title={Before you move on}]
What does \ar{معرض} mean? (\textquotedblleft{}An exhibition\textquotedblright{}.) Say it once more.
\end{tcolorbox}
\section[\texorpdfstring{\ar{نحات} (naḥḥāt)}{نحات (naḥḥāt)}]{\ar{نحات} (naḥḥāt) — a sculptor}

\label{lesson:AR-C541-nahhat}



\begin{quote}
Before the new one: say the Arabic for an artist, then the Arabic for an exhibition.
\end{quote}

\subsection*{You'll want to know: \ar{نحات}}
\textbf{\ar{نحات}} — \emph{naḥḥāt} — \textquotedblleft{}a sculptor\textquotedblright{}.

\begin{quote}
\ar{صنع} \ar{النحات} \ar{التمثال} \ar{من} \ar{الرخام}.
\end{quote}

\begin{quote}
\emph{ṣanaʿa an-naḥḥāt at-timthāl min ar-rukhām.}
\end{quote}

\begin{quote}
The sculptor made the statue out of marble.
\end{quote}

The pattern faʿʿāl names a trade.

\begin{cousinweb}
\textbf{\ar{نحات}} (\emph{naḥḥāt}) is built on the native Arabic root \emph{n-ḥ-t}, \textquotedblleft{}to hew, to carve\textquotedblright{} (Lane). Naḥḥāt is the craftsman who carves.
\end{cousinweb}

\begin{grammarlens}[title={What you've built}]
One more word for describing what you lived and felt.
\end{grammarlens}

\subsection*{Your turn}
\begin{itemize}
\raggedright
  \item \emph{Say it:} \emph{naḥḥāt}
  \item \emph{Say it:} \emph{naḥḥāt}, once more
  \item \emph{Say it:} \emph{ṣanaʿa an-naḥḥāt at-timthāl min ar-rukhām.} — The sculptor made the statue out of marble.
  \item \emph{Recall:} say the Arabic for an artist, then the Arabic for an exhibition, then say \emph{naḥḥāt} again
\end{itemize}

\begin{tcolorbox}[breakable,colback=teal!4,colframe=teal!35!black,title={Before you move on}]
What does \ar{نحات} mean? (\textquotedblleft{}A sculptor\textquotedblright{}.) Say it once more.
\end{tcolorbox}
\section[\texorpdfstring{\ar{نحت} (naḥata)}{نحت (naḥata)}]{\ar{نحت} (naḥata) — to carve}

\label{lesson:AR-C541-nahata}



\begin{quote}
Before the new one: say the Arabic for an exhibition, then the Arabic for a sculptor.
\end{quote}

\subsection*{You'll want to know: \ar{نحت}}
\textbf{\ar{نحت}} — \emph{naḥata} — \textquotedblleft{}to carve\textquotedblright{}.

\begin{quote}
\ar{نحت} \ar{الفنان} \ar{حصانا} \ar{من} \ar{الخشب}.
\end{quote}

\begin{quote}
\emph{naḥata al-fannān ḥiṣānan min al-khashab.}
\end{quote}

\begin{quote}
The artist carved a horse out of wood.
\end{quote}

Present: \ar{ينحت}, \emph{yanḥit}.

The noun \ar{نحت}, read naḥt, is sculpture.

\begin{cousinweb}
\textbf{\ar{نحت}} (\emph{naḥata}) is built on the native Arabic root \emph{n-ḥ-t}, \textquotedblleft{}to hew, to carve\textquotedblright{} (Lane). Lane gives yanḥitu as its present.
\end{cousinweb}

\begin{grammarlens}[title={What you've built}]
One more word for describing what you lived and felt.
\end{grammarlens}

\subsection*{Your turn}
\begin{itemize}
\raggedright
  \item \emph{Say it:} \emph{naḥata}
  \item \emph{Say it:} \emph{naḥata}, once more
  \item \emph{Say it:} \emph{naḥata al-fannān ḥiṣānan min al-khashab.} — The artist carved a horse out of wood.
  \item \emph{Recall:} say the Arabic for an exhibition, then the Arabic for a sculptor, then say \emph{naḥata} again
\end{itemize}

\begin{tcolorbox}[breakable,colback=teal!4,colframe=teal!35!black,title={Before you move on}]
What does \ar{نحت} mean? (\textquotedblleft{}To carve\textquotedblright{}.) Say it once more.
\end{tcolorbox}
\section[\texorpdfstring{\ar{رسام} (rassām)}{رسام (rassām)}]{\ar{رسام} (rassām) — a painter (artist)}

\label{lesson:AR-C541-rassam}



\begin{quote}
Before the new one: say the Arabic for a sculptor, then the Arabic for to carve.
\end{quote}

\subsection*{You'll want to know: \ar{رسام}}
\textbf{\ar{رسام}} — \emph{rassām} — \textquotedblleft{}a painter (artist)\textquotedblright{}.

\begin{quote}
\ar{رسم} \ar{الرسام} \ar{البحر} \ar{عند} \ar{الغروب}.
\end{quote}

\begin{quote}
\emph{rasama ar-rassām al-baḥr ʿinda al-ghurūb.}
\end{quote}

\begin{quote}
The painter painted the sea at sunset.
\end{quote}

The pattern faʿʿāl names a trade or a habit.

\begin{cousinweb}
\textbf{\ar{رسام}} (\emph{rassām}) is built on the native Arabic root \emph{r-s-m}, \textquotedblleft{}to mark, to trace\textquotedblright{} (Lane). Rasama, to draw, which you know, is from the same root.
\end{cousinweb}

\begin{grammarlens}[title={What you've built}]
One more word for describing what you lived and felt.
\end{grammarlens}

\subsection*{Your turn}
\begin{itemize}
\raggedright
  \item \emph{Say it:} \emph{rassām}
  \item \emph{Say it:} \emph{rassām}, once more
  \item \emph{Say it:} \emph{rasama ar-rassām al-baḥr ʿinda al-ghurūb.} — The painter painted the sea at sunset.
  \item \emph{Recall:} say the Arabic for a sculptor, then the Arabic for to carve, then say \emph{rassām} again
\end{itemize}

\begin{tcolorbox}[breakable,colback=teal!4,colframe=teal!35!black,title={Before you move on}]
What does \ar{رسام} mean? (\textquotedblleft{}A painter (artist)\textquotedblright{}.) Say it once more.
\end{tcolorbox}
